\begin{resumo}[Abstract]
 \begin{otherlanguage*}{english}
   This work presents the reconstruction of the web platform of Pró-Mamá, a municipal breastfeeding program in Osório, from the perspective of technological continuity and digital support for maternal and child health. The research is characterized as applied, technological in nature and based on a mixed approach, organized as a single case study and supported by documentary research, requirements analysis and incremental development. The work reviewed limitations of the previous version, especially the absence of secure communication, maintenance difficulties, low infrastructure reproducibility and dependence on fixed rules embedded in the source code. As a result, the API, the Administrative Panel and the deployment infrastructure were restructured using technologies such as Node.js, TypeScript, React.js, PostgreSQL, Redis, Docker, Nginx and Cloudflare. The solution now includes profile-based authentication and authorization, automated OpenAPI documentation, asynchronous notification processing, controlled migration of legacy data, HTTPS publication and continuous integration and delivery routines. Automated validation focused on the back-end, totaling 143 successful checks and global coverage close to 80\% for line and statement metrics. It is concluded that the reconstruction delivered a safer, more manageable and sustainable foundation for the continuity of Pró-Mamá, preserving its social value and expanding its capacity for maintenance and future growth.

   \vspace{\onelineskip}
 
   \noindent 
   \textbf{Keywords}: breastfeeding. Pró-Mamá. digital health. web applications. information security.
 \end{otherlanguage*}
\end{resumo}
